\section{Low-rank cross-observable covariance theorem}
\label{sec:bhsm-optical-low-rank}

Even before the compact full-preimage spatial-current background is
numerically available, the optical construction makes a normalization-
independent prediction.  Let $q\in\mathbb R^m$ contain the active physical
topographic texture amplitudes and let $M_O$ be the complete linear response
from those amplitudes to a fixed observable vector $O$.  Then
\begin{equation}
\boxed{
C_O^{\rm topo}
=
M_O C_q M_O^\dagger
}
\end{equation}
and hence
\begin{equation}
\boxed{
{\rm rank}\,C_O^{\rm topo}
\le
{\rm rank}\,C_q
\le
m.
}
\end{equation}

For one active texture channel,
\begin{equation}
C_O^{\rm topo}
=
\sigma_q^2 m_Om_O^\dagger,
\end{equation}
so every pair satisfies
\begin{equation}
\boxed{
C_{ii}C_{jj}-|C_{ij}|^2=0.
}
\end{equation}
Equivalently, in the noiseless isolated limit,
\begin{equation}
R_1=\frac{\lambda_1}{\sum_i\lambda_i}=1,
\qquad
\rho_{21}=\frac{\lambda_2}{\lambda_1}=0.
\end{equation}

The rank statement applies to the BHSM topographic contribution, not to raw
survey covariance.  The test object is
\begin{equation}
\widehat C_{\rm res}
=
\widehat C_{\rm obs}
-
\widehat C_{\rm std}
-
\widehat N,
\end{equation}
with the standard-physics and measurement covariances constructed
independently of the BHSM rank target.

A residual covariance requiring more than the action-allowed number of
independent eigenmodes falsifies that BHSM texture-channel count even if
individual pairwise correlations are nonzero.  This test does not require the
absolute seam-export normalization.
